% !TEX program = xelatex
\documentclass[12pt,a4paper]{article}

\usepackage{fontspec}
\setmainfont{Arial} % FAPESP APR: fonte Calibri ou Arial, tamanho 12
\usepackage[brazilian]{babel}
\usepackage{microtype}
\usepackage[a4paper,left=3cm,right=1.5cm,top=2.5cm,bottom=2.5cm]{geometry} % FAPESP APR: margens 3cm esq. / 1.5cm dir.
\usepackage{setspace}
\onehalfspacing % FAPESP APR: espaçamento entre linhas de 1.5
\usepackage{amsmath,amssymb}
\usepackage{booktabs}
\usepackage{longtable}
\usepackage{array}
\usepackage{enumitem}
\usepackage{xcolor}
\usepackage{titlesec}
\usepackage[numbers,sort&compress]{natbib}
\usepackage[hidelinks,breaklinks]{hyperref}

\setlist{itemsep=2pt,topsep=4pt}
\titleformat{\section}{\normalfont\large\bfseries}{\thesection}{0.75em}{}
\titleformat{\subsection}{\normalfont\normalsize\bfseries}{\thesubsection}{0.75em}{}

\newcommand{\pamr}{P\nobreakdash-AMR}

\title{\textbf{Precision-AMR: Precisão de Ponto Flutuante Espacialmente
Adaptativa para Migração Reversa no Tempo Baseada em EDPs}\\[0.5em]
\large\normalfont\itshape Um novo eixo de adaptividade numérica para
solucionadores adjuntos da equação da onda, desenvolvido em hardware NVIDIA
dedicado e validado em arquiteturas heterogêneas de GPUs NVIDIA e AMD no
supercomputador Santos Dumont (LNCC)}

\author{}
\date{}

\begin{document}

\maketitle
\thispagestyle{empty}

\begin{flushleft}
\textbf{Pesquisador Principal:} Prof.\ Sandro Rigo\\
\textbf{Cargo/Função:} Professor Associado, Instituto de Computação\\
\textbf{Universidade:} Universidade Estadual de Campinas (UNICAMP)\\[0.4em]
\textbf{Colaborador 1:} Thiago Maltempi --- Pesquisador de Pós-Graduação
(Doutorando), Instituto de Computação, Universidade Estadual de Campinas\\
\textbf{Colaborador 2:} Prof.\ Hervé Yviquel --- Professor Assistente,
Instituto de Computação, Universidade Estadual de Campinas
\end{flushleft}

\section{Enunciado do Problema}

O hardware de GPU vem se voltando cada vez mais para a baixa precisão em prol
de cargas de trabalho de IA (FP8/BF16), expondo novas hierarquias de memória,
tensor cores e bibliotecas matemáticas de precisão mista. Ao mesmo tempo,
simulações científicas baseadas em equações diferenciais parciais (EDPs)
continuam a exigir precisão FP64/FP32 para permanecer confiáveis. A Migração
Reversa no Tempo (Reverse Time Migration, RTM), método de imageamento sísmico
que resolve a equação da onda por diferenças finitas ao longo de milhares de
passos de tempo, é um exemplo canônico dessa tensão: geofísicos usam a imagem
migrada para orientar decisões de perfuração de alto custo, de modo que a
fidelidade numérica é essencial. A RTM também é intensiva em memória,
exigindo centenas de gigabytes de estado do campo de onda e dependendo de
checkpointing para trocar memória por recomputação em sua estrutura adjunta
de propagação direta/reversa.

Nosso trabalho anterior atacou o principal gargalo dessa estrutura, isto é, a
transferência de checkpoints entre host e GPU via PCIe, com a combinação do conceito de prefetching com a compressão de dados.
Essa ideia foi materializada como a biblioteca GPUZIP~\cite{maltempi2025checkpointing,gpuzip}, que comprime os checkpoints
na GPU e os pré-busca de forma proativa, reduzindo o volume de transferência
host--GPU em até 80$\times$ e alcançando ganhos de desempenho ponta a ponta
de até 5,1$\times$ na V100, preservando a qualidade da imagem (SSIM
$\approx$ 1, erro relativo próximo de zero). Separadamente, já se demonstrou
que a aritmética de precisão reduzida é viável para modelagem sísmica, RTM e
FWI: Fabien-Ouellet~\cite{fabienouellet2020half} constatou que uma
implementação estável em meia precisão (FP16) da equação de onda elástica é
alcançável por meio de escalonamento (\emph{scaling}) apropriado, que o erro
resultante no sismograma é uma fração desprezível da energia sísmica total e
majoritariamente incoerente com as fases sísmicas, e que esse ruído não
prejudica a qualidade das imagens de FWI ou RTM. Entretanto, essa redução é
aplicada de forma uniforme em todo o domínio computacional. Gao e
Harms~\cite{gao2023half} identificaram um risco específico dessa abordagem:
a integração no tempo em meia precisão de forma ingênua acumula erro de
arredondamento por meio de uma somatória recursiva disfarçada na regra de
atualização, degradando a qualidade da solução, sendo a somatória compensada
(de Kahan) um remédio eficaz.

Ao mesmo tempo, o refinamento adaptativo de malha (Adaptive Mesh Refinement,
AMR) tem sido explorado na simulação de ondas sísmicas e, mais recentemente,
explicitamente em fluxos de trabalho de FWI e RTM, para evitar o desperdício
inerente a uma malha uniforme: os erros de dispersão e dissipação exigem
espaçamento de malha proporcional ao menor comprimento de onda, de modo que
uma malha uniforme é forçada a usar o espaçamento mais fino em toda parte,
mesmo onde o campo de velocidade local toleraria um espaçamento bem mais
grosseiro. O AMR resolve isso refinando o $\Delta x$ localmente. Não existe,
porém, esquema comparável e fisicamente motivado para a precisão de ponto
flutuante: os trabalhos existentes de precisão mista em imageamento sísmico
são globais e uniformes (como acima), ou adaptativos apenas no nível de uma
operação individual de álgebra linear (GEMM de precisão mista, refinamento
iterativo, emulação pelo esquema de Ozaki), sem estar vinculados a uma
decomposição espacial do próprio campo de onda. Nenhum desses trabalhos varia
a precisão de ponto flutuante segundo critérios espaciais fisicamente
motivados. Esse é um terceiro eixo de adaptabilidade que permanece, portanto,
praticamente inexplorado.

Este projeto propõe preencher essa lacuna com o Precision-AMR (\pamr{}): um
esquema espacialmente estruturado em blocos, análogo de precisão adaptativa
ao AMR, que mantém fixa a resolução da malha e, em vez disso, varia o formato
de ponto flutuante atribuído a cada bloco (\emph{tile}) do domínio, com base
em critérios de tolerância a erro fisicamente motivados (proximidade a bordas
absorventes, densidade de energia do campo de onda, magnitude da
sensibilidade adjunta, e uma faixa dependente do tempo que segue a frente de
onda em propagação) que preveem se o erro de arredondamento introduzido
naquele bloco pode alcançar a condição de imageamento ou o gradiente do FWI.
Trata-se de um problema estruturalmente diferente do AMR baseado em malha. No
AMR, o artefato característico é uma reflexão localizada na interface
grossa-fina da malha, porque o refinamento altera a própria discretização.
No \pamr{}, nenhuma fronteira de discretização é introduzida; em vez disso,
um nível espacialmente variável de ruído de arredondamento é injetado em uma
EDP hiperbólica, e esse ruído se propaga junto com a onda. Os critérios de
marcação (\emph{tagging}), portanto, não podem se basear apenas na suavidade
local, como fazem os indicadores de refinamento do AMR; eles precisam, em vez
disso, aproximar onde um erro, uma vez injetado, decairá, será absorvido ou
deixará de alcançar os receptores ou o kernel de sensibilidade antes do fim
da simulação, e controlar o acúmulo de arredondamento ao longo dos milhares
de passos de tempo típicos de RTM/FWI.

Como a RTM/FWI já exige uma decomposição do domínio em blocos/\emph{tiles}
para o checkpointing, o \pamr{} reaproveita essa mesma decomposição como
granularidade de marcação, unificando o orçamento de precisão e o orçamento
de compressão com perdas do checkpointing diferencial em uma única meta de
precisão por bloco, conjuntamente ajustada. O GEMM de precisão mista, o
refinamento iterativo do cuSOLVER e a emulação de FP64 pelo esquema de Ozaki
continuam essenciais como técnicas de suporte, mas passam a ser aplicados
seletivamente, apenas onde os critérios do \pamr{} indicam necessidade ---
a forma natural de explorar GPUs orientadas primariamente para IA, como a RTX
PRO 6000 Blackwell, que oferecem vazão nativa de FP64 desprezível.

Como os critérios de marcação são motivados pela física da propagação de
ondas e pela sensibilidade adjunta, e não pelas unidades aritméticas de um
fabricante específico, tratamos sua portabilidade entre arquiteturas como uma
questão explícita, e não como uma suposição. Nosso grupo de pesquisa tem acesso ao supercomputador
Santos Dumont (LNCC/SINAPAD), que inclui aceleradores AMD
Instinct MI300A ao lado de diversas gerações NVIDIA (Volta V100, Hopper H100,
Grace-Hopper GH200). Isso nos permite testar essa questão diretamente, em vez
de apenas afirmá-la, e é o que confere ao Precision-AMR sua contribuição
potencial para a área: um eixo de adaptividade numérica novo e portável entre
fabricantes, e não apenas uma otimização pontual para uma única arquitetura
de GPU.

\section{Resultados Esperados}

A principal contribuição deste projeto é o próprio
método Precision-AMR (\pamr{}), que explora um novo eixo de adaptividade numérica
que varia a precisão de ponto flutuante espacialmente, segundo critérios
fisicamente motivados (proximidade à PML, densidade de energia do campo de
onda, magnitude da sensibilidade adjunta e faixa de esteira da frente de
onda), proposto e validado especificamente para solucionadores adjuntos de
EDPs como RTM/FWI. Essa contribuição inclui a formulação dos quatro
critérios de marcação, a caracterização do artefato de ``interface de
precisão'', análogo às reflexões de interface grossa-fina do AMR clássico, e a demonstração de que os orçamentos de precisão e de compressão com
perdas podem ser unificados em uma única meta de erro por bloco.

Como contribuições secundárias, o projeto produzirá: (i) uma avaliação
multi-arquitetura do uso de precisão mista guiado pelo P-AMR, comparando GPUs
NVIDIA (V100, H100, GH200) e a APU AMD Instinct MI300A, caracterizando
quanto do ganho de desempenho e do compromisso de precisão alcançados é
independente de fabricante; e (ii) a implementação do P-AMR como uma
\textbf{biblioteca independente}, com uma interface que permite sua
integração a outras aplicações de simulação baseadas em EDPs, para além do
escopo específico da RTM/FWI. Essa biblioteca poderá ser adotada
isoladamente ou combinada ao GPUZIP, mesclando o zoneamento espacial de
precisão do P-AMR com o checkpointing comprimido do GPUZIP em um único fluxo
de trabalho, sem exigir que a aplicação hospedeira adote toda a pilha do
GPUZIP apenas para se beneficiar do P-AMR. Ela será disponibilizada
em código aberto, sob licença MIT, no GitHub.

O resultado científico central será submetido a um evento ou periódico de
alto impacto em computação de alto desempenho, por exemplo,
Supercomputing (SC), IPDPS, Euro-Par ou SBAC-PAD entre os eventos, ou a
IJHPCA ou um periódico do tipo \emph{Transactions} da IEEE ou da ACM entre
os periódicos, acompanhado de um pacote público de reprodutibilidade
(conjuntos de dados, perfis do Nsight, métricas de qualidade),
disponibilizado via REDU da UNICAMP.

% ============================================================================
% BLOCO ANTIGO (pré-reestruturação FAPESP) --- preservado para referência.
% Remover apenas após aprovação final da versão reestruturada.
% ============================================================================
\iffalse

\section*{Resumo}

Simulações científicas como a Migração Reversa no Tempo (Reverse Time
Migration, RTM) resolvem equações diferenciais parciais (EDPs) que exigem
precisão FP64/FP32, mas as arquiteturas modernas de GPU vêm sendo cada vez mais
otimizadas para cargas de trabalho de IA em baixa precisão (FP8/BF16),
expondo novas hierarquias de memória, tensor cores e bibliotecas matemáticas de
precisão mista. Trabalhos anteriores sobre aceleração de RTM reduziram a
precisão de ponto flutuante de forma uniforme em todo o domínio computacional;
separadamente, o refinamento adaptativo de malha (Adaptive Mesh Refinement,
AMR) tem sido explorado na modelagem de ondas sísmicas para variar o
espaçamento da malha de acordo com o comprimento de onda local. Este projeto
propõe um terceiro eixo de adaptividade, ainda pouco explorado: o
Precision-AMR (\pamr{}), um esquema espacialmente estruturado em blocos que
varia a precisão de ponto flutuante --- não a resolução da malha --- segundo
critérios de tolerância a erro por bloco fisicamente motivados (proximidade a
bordas absorventes, densidade de energia do campo de onda, magnitude da
sensibilidade adjunta e uma faixa dependente do tempo que segue a frente de
onda em propagação).

Diferentemente do AMR baseado em malha, em que os artefatos numéricos ficam
localizados nas interfaces grossas-finas da malha, o ruído de arredondamento
induzido pela precisão se propaga junto com a própria onda; assim, os
critérios de marcação (tagging) do \pamr{} devem identificar onde o erro
injetado não alcançará a condição de imageamento nem o gradiente do FWI, além
de controlar o acúmulo de arredondamento ao longo dos milhares de passos de
tempo típicos de RTM/FWI. Construímos o \pamr{} sobre nosso framework GPUZIP
anterior, que acelerou a RTM em sistemas V100 por meio de compressão com
perdas (\emph{lossy compression}) realizada na GPU (nvCOMP/Bitcomp) e
pré-busca (\emph{prefetching}) de checkpoints, reduzindo o volume de
transferência host--GPU em até 80$\times$ e proporcionando um ganho de
desempenho ponta a ponta de até 5,1$\times$. O checkpointing diferencial, a
álgebra linear de precisão mista via cuBLAS/cuSOLVER, a emulação de FP64 pelo
esquema de Ozaki e os operadores espectrais baseados em cuFFT são mantidos
como técnicas de suporte, agora aplicadas seletivamente, por bloco (tile), de
acordo com os critérios estabelecidos pelo \pamr{}, em vez de uniformemente em
todo o domínio.

O programa de 12 meses combina dois recursos computacionais. Uma GPU RTX PRO
6000 Blackwell, solicitada como hardware local dedicado, sustenta o
desenvolvimento cotidiano e o estudo de emulação seletiva central à Fase 3,
já que sua vazão nativa de FP64 desprezível torna a questão do zoneamento de
precisão inevitável, e não apenas opcional. O supercomputador Santos Dumont
(LNCC/SINAPAD) --- cujas partições abrangem aceleradores NVIDIA Volta V100,
Hopper H100 e Grace-Hopper GH200, além de APUs AMD Instinct MI300A (CDNA3)
--- sustenta a varredura de validação em larga escala entre conjuntos de
dados, algoritmos de checkpointing e critérios de marcação do \pamr{}, e nos
permite testar diretamente se a metodologia de marcação generaliza para além
de um único fabricante de GPU. Os resultados esperados incluem uma versão de
código aberto do GPUZIP v3.0 construída em torno de um módulo Precision-AMR, e
uma publicação revisada por pares.

\paragraph{Palavras-chave:} migração reversa no tempo; precisão adaptativa;
computação de precisão mista; aritmética de ponto flutuante; solucionadores
adjuntos de EDPs; adaptividade estruturada em blocos; portabilidade entre
arquiteturas; computação em GPU; compressão com perdas; checkpointing.

\fi
% ============================================================================
% FIM DO BLOCO ANTIGO (Resumo)
% ============================================================================

\section{Plataformas Computacionais}

\paragraph{HPC SDK} (nvc++/nvfortran, bibliotecas matemáticas CUDA-X, Nsight
Systems/Compute). Experiência: a equipe usa CUDA e Nsight Systems
rotineiramente. Adoção: migrar a compilação do Awave-3D de clang para nvc++ e
validar a paridade numérica em relação às linhas de base (\emph{baselines})
existentes. O Nsight Compute fornece análise \emph{roofline} em nível de
kernel, agora aplicada por bloco espacial, para identificar quais blocos são
limitados pela largura de banda de memória versus limitados por computação, e
para delimitar objetivamente onde um bloco de menor precisão realmente
compensa.

\paragraph{RTX PRO 6000 Blackwell} (hardware dedicado, solicitado). Os tensor
cores de 5ª geração dessa GPU, com suporte a FP8/FP4 e vazão nativa de FP64
desprezível, a tornam a plataforma central para o \pamr{}: como a emulação de
FP64 em todo o domínio seria cara em toda parte, a questão científica que este
projeto responde é precisamente quais blocos precisam dela e quais não
precisam. Sua memória de 96\,GB GDDR7 com ECC também sustenta o
desenvolvimento local; por se tratar de hardware dedicado e sem fila de
espera, permite depuração iterativa rápida, algo para o qual um sistema de
lote compartilhado não é bem adequado.

\paragraph{Santos Dumont} (LNCC/SINAPAD, supercomputador nacional Tier-0,
Petrópolis-RJ). Fornece o ambiente de validação em larga escala. Sua partição
BullSequana X1000 inclui 94 nós com 4$\times$ GPUs NVIDIA Volta V100 cada ---
compatível com o hardware usado nas linhas de base originais do GPUZIP ---
além de um nó de IA dedicado com 8$\times$ V100 de 16GB via NVLink. Sua
partição mais recente, BullSequana XH3000, acrescenta 62 nós com 4$\times$
NVIDIA Hopper H100 SXM de 80GB via NVLink, 36 nós com superchips NVIDIA
Grace-Hopper GH200, e 18 nós com 2$\times$ APUs AMD Instinct MI300A (núcleos
de GPU CDNA3, 128\,GB de HBM3 cada). Essa heterogeneidade nos permite gerar a
referência de verdade (\emph{ground truth}) em FP64 e executar em escala a
varredura completa dos critérios do \pamr{} em hardware NVIDIA, além de testar
separadamente se a metodologia de marcação e a abordagem de tiling em blocos
generalizam para uma arquitetura não NVIDIA.

\paragraph{nvCOMP / Bitcomp} (compressor com perdas executado na GPU).
Experiência: integrado e validado no GPUZIP v2.0, onde o Bitcomp apresentou o
melhor equilíbrio entre precisão e compressão na V100. Adoção aqui: o
mecanismo de armazenamento do checkpointing diferencial, agora indexado pelos
mesmos identificadores de bloco usados na marcação de precisão, de modo que
cada bloco carrega um único orçamento de erro conjuntamente ajustado para
precisão e compressão.

\paragraph{cuBLAS (extensões de precisão mista) e cuSOLVER (refinamento
iterativo).} Ambos vêm incluídos no HPC SDK. Experiência: a equipe já utilizou
rotinas BLAS/LAPACK em trabalhos anteriores de HPC. Adoção: aplicados apenas
nos blocos que o \pamr{} sinaliza como contendo subkernels densos que
necessitam de precisão além da baixa precisão nativa (computação em FP32 com
acumulação em FP64, e refinamento iterativo do cuSOLVER).

\paragraph{Emulação de FP64 pelo esquema de Ozaki} via tensor cores de baixa
precisão. Adoção: reservada ao subconjunto de blocos que os critérios de
sensibilidade do \pamr{} marcam como críticos para a precisão (próximos à
frente de onda, com alta magnitude de sensibilidade adjunta), em vez de
aplicada uniformemente --- testando diretamente se a emulação seletiva é mais
barata do que emular em todo o domínio, preservando a fidelidade da imagem.
Trabalhos recentes demonstraram ganhos de até 2,3$\times$ em relação ao DGEMM
nativo em FP64 na RTX PRO 6000 Blackwell~\cite{uchino2025ozaki}.

\paragraph{cuFFT.} Experiência: utilizado em trabalhos anteriores com
operadores espectrais. Adoção: avaliado como alternativa aos estênceis de
diferenças finitas (FDM) para a atualização do campo de onda dentro de blocos
selecionados, marcados por precisão, caracterizando o compromisso entre
largura de banda de memória e computação em relação ao caminho padrão de
diferenças finitas.

\paragraph{AMD Instinct MI300A} (CDNA3, Santos Dumont, exploratório). Como
atividade adicional de validação, avaliamos quanto do framework \pamr{} --- a
decomposição em blocos/tiles, os quatro critérios de marcação e o orçamento
conjunto de erro de precisão e compressão --- se transfere para o ecossistema
ROCm/HIP da AMD (rocBLAS, hipBLAS e os próprios caminhos de matriz de precisão
reduzida da AMD), sem assumir paridade de recursos com a emulação do esquema
de Ozaki, específica da CUDA, ou com o caminho de compressão nvCOMP/Bitcomp. O
objetivo é determinar se o Precision-AMR é uma metodologia portável ou uma
otimização específica da NVIDIA.

\section{Conjunto de Dados e Modelo}

\begin{itemize}
  \item Tamanho: $\approx$X\,TB no total (a confirmar junto ao REDU). Os
        conjuntos de checkpoints chegaram a $\approx$330\,GB de RAM do host na
        V100; um único checkpoint do Marmousi3D tem $\approx$500\,MB.
  \item Tipo de dado: modelos de velocidade sísmica e traços (malhas
        numéricas).
  \item Origem: Salt (SEG/EAGE) e Marmousi3D são conjuntos de dados abertos; o
        conjunto de dados Large é interno (\emph{in-house}).
  \item Confidencialidade: não confidencial, sem dados pessoais, sem NDA.
  \item Parâmetros de modelo de IA: não aplicável (solucionador numérico de
        EDPs).
  \item Local de armazenamento: Universidade Estadual de Campinas (UNICAMP),
        Campinas, Brasil; preparado para envio aos sistemas de arquivos
        paralelos do Santos Dumont para as execuções em larga escala.
\end{itemize}

% ============================================================================
% BLOCO ANTIGO (pré-reestruturação FAPESP) --- preservado para referência.
% Conteúdo já incorporado à seção "Enunciado do Problema" acima.
% Remover apenas após aprovação final da versão reestruturada.
% ============================================================================
\iffalse

\section{Introdução}

O hardware vem se voltando para a baixa precisão em prol da IA, mas a
simulação científica baseada em EDPs exige alta precisão para permanecer
confiável. A RTM, método de imageamento sísmico que resolve a equação da onda
por diferenças finitas ao longo de milhares de passos de tempo, é um exemplo
canônico: geofísicos usam a imagem migrada para orientar decisões de
perfuração de alto custo, de modo que a fidelidade numérica é inegociável. A
RTM também é intensiva em memória, exigindo centenas de gigabytes de estado do
campo de onda e dependendo de checkpointing para trocar memória por
recomputação em sua estrutura adjunta de propagação direta/reversa.

Nosso trabalho anterior atacou o principal gargalo --- a transferência de
checkpoints entre host e GPU via PCIe --- com o
GPUZIP~\cite{maltempi2025checkpointing,gpuzip}, que comprime os checkpoints na
GPU e os pré-busca de forma proativa, alcançando ganhos de até 5,1$\times$ na
V100 preservando a qualidade da imagem (SSIM $\approx$ 1, erro relativo
próximo de zero). Separadamente, já se demonstrou que a aritmética de precisão
reduzida é viável para modelagem sísmica, RTM e FWI:
Fabien-Ouellet~\cite{fabienouellet2020half} constatou que uma implementação
estável em meia precisão (FP16) da equação de onda elástica é alcançável por
meio de escalonamento (\emph{scaling}) apropriado, que o erro resultante no
sismograma é uma fração desprezível da energia sísmica total e majoritariamente
incoerente com as fases sísmicas, e que esse ruído não prejudica a qualidade
das imagens de FWI ou RTM. Entretanto, essa redução é aplicada de forma
uniforme em todo o domínio. Gao e Harms~\cite{gao2023half} identificaram um
risco específico dessa abordagem: a integração no tempo em meia precisão de
forma ingênua acumula erro de arredondamento por meio de uma somatória
recursiva disfarçada na regra de atualização, degradando a qualidade da
solução, sendo a somatória compensada (de Kahan) um remédio eficaz.

Ao mesmo tempo, o refinamento adaptativo de malha (AMR) tem sido explorado na
simulação de ondas sísmicas e, mais recentemente, explicitamente em fluxos de
trabalho de FWI e RTM, para evitar o desperdício inerente a uma malha
uniforme: os erros de dispersão e dissipação exigem espaçamento de malha
proporcional ao menor comprimento de onda, de modo que uma malha uniforme é
forçada a usar o espaçamento mais fino em toda parte, mesmo onde o campo de
velocidade local toleraria um espaçamento bem mais grosseiro. O AMR resolve
isso refinando o $\Delta x$ localmente. Não existe esquema comparável,
fisicamente motivado, para a precisão de ponto flutuante: os trabalhos
existentes de precisão mista em imageamento sísmico são globais e uniformes
(como acima), ou adaptativos apenas no nível de uma operação individual de
álgebra linear (GEMM de precisão mista, refinamento iterativo, emulação pelo
esquema de Ozaki), sem estar vinculados a uma decomposição espacial do próprio
campo de onda.

Este projeto propõe preencher essa lacuna com o Precision-AMR (\pamr{}): um
análogo de precisão adaptativa ao AMR que mantém fixa a resolução da malha e,
em vez disso, varia o formato de ponto flutuante atribuído a cada bloco
(\emph{tile}) do domínio, com base em critérios que preveem se o erro de
arredondamento introduzido naquele bloco pode alcançar a condição de
imageamento ou o gradiente do FWI. Trata-se de um problema estruturalmente
diferente do AMR baseado em malha. No AMR, o artefato característico é uma
reflexão localizada na interface grossa-fina da malha, porque o refinamento
altera a própria discretização. No \pamr{}, nenhuma fronteira de discretização
é introduzida; em vez disso, um nível espacialmente variável de ruído de
arredondamento é injetado em uma EDP hiperbólica, e esse ruído se propaga
junto com a onda. Os critérios de marcação, portanto, não podem se basear
apenas na suavidade local (como fazem os indicadores de refinamento do AMR);
eles precisam, em vez disso, aproximar onde um erro, uma vez injetado,
decairá, será absorvido ou deixará de alcançar os receptores ou o kernel de
sensibilidade antes do fim da simulação.

Propomos quatro critérios candidatos, detalhados em Métodos: proximidade à
borda absorvente (PML), onde a energia já está sendo descartada
propositalmente; densidade de energia local do campo de onda, que rege quanto
o erro de um dado bloco pode contribuir para a correlação final; magnitude da
sensibilidade adjunta, reaproveitando o cálculo de gradiente que o FWI já
realiza; e uma faixa dependente do tempo que segue a frente de onda em
propagação, análoga a um detector de choque móvel em códigos de AMR
astrofísicos ou de combustão. Como a RTM/FWI já exige uma decomposição do
domínio em blocos/\emph{tiles} para o checkpointing, o \pamr{} reaproveita essa
mesma decomposição como granularidade de marcação, unificando o orçamento de
precisão e o orçamento de compressão com perdas do checkpointing diferencial
em uma única meta de precisão por bloco, conjuntamente ajustada. O GEMM de
precisão mista, o refinamento iterativo do cuSOLVER e a emulação pelo esquema
de Ozaki continuam essenciais, mas passam a ser aplicados seletivamente,
apenas onde os critérios do \pamr{} indicam necessidade, o que constitui a
forma natural de explorar GPUs como a RTX PRO 6000 Blackwell, que oferecem
vazão nativa de FP64 desprezível.

Como os critérios de marcação são motivados pela física da propagação de ondas
e pela sensibilidade adjunta, e não pelas unidades aritméticas de um
fabricante específico, tratamos sua portabilidade como uma questão explícita,
e não como uma suposição. O acesso ao supercomputador Santos Dumont --- cujas
partições incluem aceleradores AMD Instinct MI300A ao lado de diversas gerações
NVIDIA --- nos permite testar essa questão diretamente, em vez de apenas
afirmá-la.

\fi
% ============================================================================
% FIM DO BLOCO ANTIGO (Introdução)
% ============================================================================

\section{Métodos}

\subsection{Fase 1 --- Infraestrutura e Linha de Base (M1--M2)}

Portar o Awave-3D e o GPUZIP para o HPC SDK (nvc++), reintegrar o
nvCOMP/Bitcomp e reproduzir as linhas de base em V100 na partição V100 do
Santos Dumont, compatível com o hardware original do GPUZIP.

\begin{itemize}
  \item Estabelecer uma decomposição do domínio em blocos/\emph{tiles},
        estendendo a estrutura de blocos de checkpoint já existente para que
        também possa carregar uma marca (\emph{tag}) de precisão por bloco.
  \item Perfilar (\emph{profile}) todos os kernels com Nsight Systems e Nsight
        Compute, tanto na RTX PRO 6000 quanto no Santos Dumont, confirmando
        que os estênceis de diferenças finitas (FDM) são limitados pela
        largura de banda de memória na granularidade de bloco antes de
        qualquer alteração de precisão ser tentada.
\end{itemize}

\subsection{Fase 2 --- Precision-AMR: Critérios, Marcação e Validação
(M2--M6)}

Esta é a fase científica central, executada em escala na partição H100 do
Santos Dumont. Projetamos, implementamos e validamos quatro critérios
candidatos de marcação para o zoneamento espacial de precisão:

\begin{itemize}
  \item \textbf{Critério de borda:} blocos dentro ou adjacentes à camada
        absorvente PML são marcados por padrão para precisão reduzida, já que
        sua contribuição já está sendo amortecida.
  \item \textbf{Critério de energia:} os blocos são marcados por um limiar
        sobre a densidade de energia local do campo de onda, de modo que
        regiões de baixa amplitude (antes da chegada da frente de onda, ou
        após o espalhamento geométrico ter atenuado o sinal) tolerem menor
        precisão.
  \item \textbf{Critério de sensibilidade (FWI):} os blocos são marcados pela
        magnitude do kernel de sensibilidade adjunta local, reaproveitando o
        próprio cálculo de gradiente como sinal de marcação, em analogia
        direta a um estimador de erro de truncamento de Berger--Oliger no
        AMR clássico.
  \item \textbf{Critério de esteira da frente de onda:} uma faixa dependente
        do tempo, imediatamente atrás da frente de onda em propagação, é
        marcada para precisão reduzida, enquanto a borda dianteira --- onde a
        precisão de fase rege a condição de imageamento --- mantém precisão
        plena.
\end{itemize}

Para cada critério, implementamos a somatória compensada (de Kahan) na
atualização de integração no tempo dos blocos de baixa precisão, para
controlar o acúmulo de arredondamento ao longo de toda a simulação, seguindo o
mecanismo identificado por Gao e Harms~\cite{gao2023half}. Testamos
explicitamente a ocorrência de um artefato de ``interface de precisão'' ---
análogo às reflexões de interface grossa-fina do AMR baseado em malha --- por
meio de análise de resíduo nas fronteiras entre blocos de precisões
diferentes. Por fim, cada bloco recebe conjuntamente um formato de ponto
flutuante e uma política de compressão de checkpoint (\emph{snapshot}
completo versus delta comprimido), com a precisão combinada medida por PSNR,
SSIM e erro relativo médio em relação à linha de base em FP64, tanto no campo
de onda bruto quanto na imagem migrada final. Esta fase absorve o
checkpointing diferencial, que passa a ser o mecanismo de armazenamento para os
blocos marcados pelo \pamr{}, em vez de um entregável separado.

\subsection{Fase 3 --- Precision-AMR sem FP64 Nativo (M7--M10)}

Os critérios de marcação validados na Fase 2 são aplicados na RTX PRO 6000,
hardware com vazão nativa de FP64 desprezível. Por bloco, o \emph{pipeline}
seleciona entre execução nativa de baixa precisão em tensor cores, GEMM de
precisão mista via cuBLAS (computação em FP32 com acumulação em FP64) ou
refinamento iterativo do cuSOLVER para subkernels densos, e emulação de FP64
pelo esquema de Ozaki para blocos sinalizados como críticos para a precisão.

\begin{itemize}
  \item Isso operacionaliza diretamente a motivação de explorar GPUs
        voltadas primariamente para IA que carecem de FP64 nativo: a
        emulação, cara, é concentrada apenas onde os critérios do \pamr{}
        indicam necessidade.
  \item Quantificamos o compromisso entre ganho de desempenho e precisão em
        relação a uma linha de base com Ozaki uniforme (emulando FP64 em todo
        o domínio), para determinar quanto o \emph{tiling} seletivo realmente
        economiza.
  \item Todos os erros são delimitados conjuntamente (interação
        compressão $\times$ precisão) em relação à linha de base em FP64,
        usando PSNR, SSIM e erro relativo médio, nos mesmos três conjuntos de
        dados e três algoritmos de checkpointing (Revolve, zCut, Uniform)
        usados ao longo de todo o projeto.
  \item Atividade exploratória entre arquiteturas: a mesma lógica de marcação
        e \emph{tiling} é portada para os nós AMD Instinct MI300A do Santos
        Dumont via ROCm/HIP, para testar se os critérios e o \emph{framework}
        de orçamento de erro se sustentam em uma arquitetura CDNA3 sem acesso
        nativo à emulação pelo esquema de Ozaki ou ao nvCOMP/Bitcomp.
\end{itemize}

\subsection{Fase 4 --- Operadores Espectrais, Integração e Lançamento
(M10--M12)}

\begin{itemize}
  \item Avaliar operadores de derivada pseudo-espectrais baseados em cuFFT
        como regra de atualização alternativa para blocos selecionados,
        marcados por precisão, como técnica secundária que complementa a
        contribuição central do \pamr{}.
  \item Integrar o módulo Precision-AMR, o checkpointing diferencial e os
        caminhos de precisão mista em um lançamento unificado do GPUZIP v3.0.
  \item Avaliação de escalabilidade multi-nó no Santos Dumont
        ($\leq$8 GPUs).
  \item Análise final, redação e submissão, com o Precision-AMR como
        contribuição científica central e o checkpointing diferencial, a
        álgebra linear de precisão mista e os operadores espectrais
        apresentados como técnicas de suporte.
\end{itemize}

\section{Avaliação}

Todos os experimentos abrangem três conjuntos de dados (Large, Marmousi3D,
Salt), três algoritmos de checkpointing (Revolve, zCut, Uniform) e os quatro
critérios de marcação do \pamr{} (individualmente e em combinação). A
perfilagem em nível de kernel usa o Nsight Compute (\emph{roofline},
detalhamento de tráfego de memória); a perfilagem de linha do tempo usa o
Nsight Systems. A implementação é em C/C++ com CUDA sob o HPC SDK, com OpenMP
Cluster para escalabilidade multi-nó. A precisão é avaliada tanto no campo de
onda quanto na imagem final de RTM, já que um \emph{tiling} que preserve
PSNR/SSIM do campo de onda ainda pode, em princípio, deslocar a imagem migrada
se o erro for espacialmente correlacionado com refletores. Para a atividade
entre arquiteturas, as mesmas métricas de precisão são calculadas na partição
AMD MI300A do Santos Dumont e comparadas às linhas de base NVIDIA, para
caracterizar quanto do ganho de desempenho e do compromisso de precisão
alcançados independe da arquitetura.

% ============================================================================
% BLOCO ANTIGO (pré-reestruturação FAPESP) --- preservado para referência.
% O parágrafo "Resultados projetados" já foi incorporado à nova seção
% "Resultados Esperados" (acima, logo após "Enunciado do Problema"). O
% cronograma (tabela por mês/fase) será reaproveitado na futura seção
% "Cronograma de Execução"; o parágrafo "Dependências" ainda não foi
% realocado. Remover apenas após aprovação final da versão reestruturada.
% ============================================================================
\iffalse

\section{Resultados Esperados}

Cronograma (12 meses):

\begin{longtable}{@{}p{0.12\textwidth}p{0.5\textwidth}p{0.19\textwidth}p{0.09\textwidth}@{}}
\toprule
\textbf{Mês} & \textbf{Atividade} & \textbf{Hardware} & \textbf{Fase}\\
\midrule
\endfirsthead
\toprule
\textbf{Mês} & \textbf{Atividade} & \textbf{Hardware} & \textbf{Fase}\\
\midrule
\endhead
\bottomrule
\endfoot
M1--M2 & Portar Awave-3D + GPUZIP para o HPC SDK (nvc++); reproduzir linhas de
base em V100; estabelecer decomposição em blocos/\emph{tiles}; perfilagem
\emph{roofline} do Nsight Compute por bloco. & RTX PRO 6000; Santos Dumont
(V100) & 1\\
\addlinespace
M2--M5 & Projetar e implementar os critérios de marcação do Precision-AMR
(proximidade à PML, densidade de energia, sensibilidade adjunta, faixa de
esteira da frente de onda); somatória compensada; teste de artefato de
interface de precisão. & Santos Dumont (H100) & 2\\
\addlinespace
M5--M6 & Unificar o orçamento de precisão e compressão via checkpointing
diferencial; escalabilidade multi-nó; lançamento de código aberto do GPUZIP
v3.0-alpha; pacote de reprodutibilidade (REDU). & Santos Dumont (H100,
multi-nó) & 2\\
\addlinespace
M7--M10 & Aplicar o \emph{tiling} do Precision-AMR sem FP64 nativo: GEMM de
precisão mista via cuBLAS seletivo, refinamento iterativo do cuSOLVER e
emulação de FP64 pelo esquema de Ozaki apenas nos blocos críticos para a
precisão. Portabilidade exploratória do \emph{framework} de marcação para o
AMD Instinct MI300A via ROCm/HIP. & RTX PRO 6000; Santos Dumont (AMD MI300A,
exploratório) & 3\\
\addlinespace
M10--M12 & Avaliação de operadores espectrais cuFFT em blocos selecionados;
lançamento unificado do GPUZIP v3.0; análise final, redação e submissão do
artigo. & RTX PRO 6000; Santos Dumont & 4\\
\end{longtable}

Resultados projetados: um lançamento de código aberto do GPUZIP v3.0
(licença MIT, GitHub) construído em torno de um módulo Precision-AMR, com
checkpointing diferencial, caminhos de GEMM de precisão mista/esquema de
Ozaki e operadores espectrais cuFFT integrados como componentes de suporte;
uma publicação revisada por pares com o Precision-AMR como contribuição
central (alvo: SBAC-PAD / Euro-Par / IJHPCA); e um pacote público de
reprodutibilidade (conjuntos de dados, perfis do Nsight, métricas de
qualidade) via REDU da UNICAMP.

Dependências: esta solicitação de hardware não está condicionada à combinação
com outro financiamento. As execuções de validação em larga escala no
supercomputador Santos Dumont (LNCC/SINAPAD) ocorrem por meio do próprio
programa de alocação computacional do LNCC (SINAPAD), independentemente desta
solicitação. O apoio institucional existente financia pessoal; nenhuma dessas
fontes duplica a solicitação da RTX PRO 6000.

\fi
% ============================================================================
% FIM DO BLOCO ANTIGO (Resultados Esperados)
% ============================================================================

\section{Detalhes do Apoio ao Projeto}

Solicitamos uma GPU RTX PRO 6000 Blackwell como hardware de desenvolvimento
local dedicado. Sua abundante vazão de tensor cores em FP8/FP4 e sua vazão
nativa de FP64 desprezível a tornam o alvo cientificamente relevante para o
estudo de emulação seletiva central à Fase 3: ela força a questão de saber se
o uso seletivo de emulação pelo esquema de Ozaki pelo Precision-AMR consegue
igualar a emulação em todo o domínio a uma fração do custo. Por se tratar de
hardware dedicado e sem fila de espera, também sustenta o desenvolvimento
iterativo cotidiano para o qual um sistema de lote compartilhado não é bem
adequado. A validação em larga escala --- a varredura dos critérios do \pamr{}
entre conjuntos de dados e algoritmos de checkpointing, a geração da linha de
base em FP64, a escalabilidade multi-GPU e o teste de portabilidade para o AMD
MI300A --- é executada separadamente no supercomputador Santos Dumont
(LNCC/SINAPAD), acessado por meio do próprio programa de alocação do LNCC e
fora do escopo desta solicitação.

\subsection*{Solicitação de Recursos}

\begin{longtable}{@{}p{0.3\textwidth}p{0.62\textwidth}@{}}
\toprule
\textbf{Recurso} & \textbf{Solicitação}\\
\midrule
\endhead
\bottomrule
\endfoot
Hardware físico & 1$\times$ RTX PRO 6000 Blackwell --- hardware de
desenvolvimento local dedicado, utilizado ao longo de todo o projeto\\
\end{longtable}

O tempo de computação no Santos Dumont é solicitado separadamente por meio do
próprio programa de alocação de projetos do LNCC (SINAPAD) e não faz parte
desta solicitação de hardware.

\bibliographystyle{plainnat}
\bibliography{mixed-precision-amr-pt}

\end{document}
